\documentclass[12pt]{article}
\usepackage{amsmath}
%\usepackage{fullpage}
\usepackage[top=1in, bottom=1in, left=0.8in, right=1in]{geometry}
\usepackage{multicol}
\usepackage{wrapfig}
\usepackage{listings}
\usepackage{enumerate}
\usepackage[utf8]{inputenc}
\usepackage{graphicx}
\usepackage{caption}
\usepackage{subcaption}
\usepackage{hyperref}
\lstset{language=Java,
basicstyle={\small\ttfamily},
columns=flexible,
belowskip=0mm}

\setlength{\columnsep}{0.1pc}

\title{PHYS 24L Final Assignment}
\author{Your Name -- \texttt{UCSBNETID@umail.ucsb.edu} -- PERM NUMBER}
\date{DUE: Thursday, March 23, 2017}
\begin{document}

  \maketitle

  \vspace{-0.3in}
  \noindent
  \rule{\linewidth}{0.4pt}

  \noindent
  \textbf{For each problem, briefly explain/justify how you
  obtained your answer.} This will help us determine your understanding of
  the problem whether or not you got the correct answer. Moreover, in the
  event of an incorrect answer, we can still try to give you partial credit
  based on the explanation you provide. 


%%%%%%%%%%%%%%%%%%%%%%%%%%%%%%%%%%%%%%%%%%%%%%%%%%%%%%%%%%%%%%%%%%%%%%%%%%%%%%%%
% Problems:
\vspace{20 pt}
\noindent
\textit{Answer each of the following questions using the images provided.  (You can see the images on the next page, and download the actual files from the Overleaf Project at this link:}
\\

\centerline{
\url{https://www.overleaf.com/read/jzvpkpchxmkb\#/30252658/}
}
\vspace{20pt}
\noindent
\textit{Remember to specify the uncertainty in your measurement and to propagate uncertainties through your calculations.  You may find the following formula useful:}

\begin{eqnarray}
\sigma_f = \sqrt{\left(\frac{\partial f}{\partial x }\,\sigma_x\right)^2+\left(\frac{\partial f}{\partial y }\,\sigma_y\right)^2+\left(\frac{\partial f}{\partial z }\,\sigma_z\right)^2}
\end{eqnarray}
\vspace{20 pt}
  \begin{enumerate}

    % Problem Statement
     \item How long is the object?



      \item What is its outer diameter?



      \item What is its inner diameter?
      
      
      
      \item What is the volume of the object?
      
      
      
      \item Given that the density of the material it is made of (nylon) is (1.15 $\pm$ 0.05) g/cm$^3$, what is the mass of the object?


\end{enumerate}


\begin{figure}[p]
	\centerline{\includegraphics[width=0.9\textwidth]{nylon_stand-off_side_view.jpg}}
    	\centerline{\includegraphics[width=0.9\textwidth]{nylon_stand-off_top_view.jpg}}
    	%\caption{\small\label{fig:side_view}}
\end{figure}
  

\end{document}